\naslov{Nevronske mreže}

Nevronske mreže so sestavljene iz nevronov in sinaps.
Vsak nevron hrani stanje $v \in \R$ iz izhod $y$, sinapse pa so utežene povezave
med nevroni.
Uteži označimo z $w_i$.
Poleg tega ima vsak nevron aktivacijsko funkcijo $\phi$, s pomočjo katere
izračuna izhod: $y = \phi(v)$.
Pri evaluaciji stanje nevrona izračunamo iz stanj njegovih vhodnih sinaps kot
\[
  v = w_0 + \sum_i w_i x_i.
\]
Mrežo ustavimo, ko se neha spreminjati, zato običajno nima ciklov.

Pri usmerjeni nevronski mreži imamo nevrone razporejene v $L+2$ plasti.
Plasti $0$ pravimo \pojem{vhodna plast}, kjer imamo en nevron za vsako vhodno
spremenljivko, plasti $L+1$ pa \pojem{izhodna plast}.
Tu je en nevron v primeru regresije in en nevron za vsak razred v primeru
klasifikacije.
Preostane $L$ \pojem{skritih plasti}, ki nam dovoljujejo kompleksnejše
računanje.
Sinapse so postavljene tako, da je vsak nevron v neki plasti povezan z vsemi v
naslednji plasti, ni pa povezan s prejšnjimi plastmi ali svojo plastjo.

\vprasanje{Opiši postavitev usmerjene nevronske mreže.}

Najenostavnejša mreža je \pojem{enostavni perceptron}.
To je pravzaprav (skoraj) linearni model z $L = 0$ in stopničasto aktivacijsko
funkcijo $\phi(v) = \mathbbm{1}(v > 0)$.
Enostavneje ga zapišemo kot
\[
  \hat{y} = \phi(\mb{w}^T \mb{x}).
\]
Učenje poteka z gradientnim spustom, torej imamo predpis
\[
  \Delta \mb{w} = \eta (y - \phi(\mb{w}^T \mb{x})) \mb{x},
\]
kjer v vsaki iteraciji obravnavamo en primer, parameter $\eta$ pa se imenuje
\pojem{learning rate}.

\vprasanje{Opiši enostavni perceptron.}

Za splošnejšo klasifikacijsko nevronsko mrežo imamo v izhodnem sloju toliko
nevronov, kolikor je možnosti v klasifikaciji.
Kot odgovor mreže dobimo neka števila v $\R$, ki jih pretvorimo v verjetnosti s
funkcijo \pojem{softmax}
\[
  \phi(v_i) = \frac{e^{v_i}}{\sum_j e^{v_j}}.
\]

\vprasanje{Kaj je softmax?}

Če imamo več plasti, za gradientni spust potrebujemo odvode napake po
$w_{ji}^{l}$, torej po utežeh med $l$-to in $(l+1)$-to plastjo.
Pri tem odvajamo posredno:
\[
  \frac{\partial E}{\partial w_{ji}^l}
  = \frac{\partial E}{\partial y_i^l} \frac{\partial y_i^l}{\partial v_i^l}
  \frac{\partial v_i^l}{\partial w_{ji}^l}
  = \frac{\partial E}{\partial y_i^l} \phi'(v_i^l) y_j^{l-1},
\]
odvod napake po $y_i^l$ pa izračunamo kot
\[
  \frac{\partial E}{\partial y_i^l}
  = \sum_k \frac{\partial E}{\partial v_k^{l+1}} \frac{\partial
	v_k^{l+1}}{\partial y_i^l}
  = \sum_k \frac{\partial E}{\partial y_k^{l+1}} \frac{\partial
	y_k^{l+1}}{\partial v_k^{l+1}} \frac{\partial v_k^{l+1}}{\partial y_i^l}
  = \sum_k \frac{\partial E}{\partial y_k^{l+1}} \phi'(v_k^{l+1}) \frac{\partial
	v_k^{l+1}}{\partial y_i^l},
\]
kjer seštevamo po nevronih v naslednji plasti.
Računamo lahko torej od desne proti levi, da dobimo vse odvode.
V plasti $L+1$ namesto drugega predpisa za $\partial E / \partial y_i^{L+1}$
uporabimo $-(y_i - y_i^{L+1})$, ki izhaja iz predpisa napake $E = \pol \sum_i
(y_i - \hat{y}_i)^2$ in dejstva, da dejansko odvajamo po $\hat{y}_i$.

\vprasanje{Izpelji predpis za računanje gradienta v večplastni nevronski mreži.}

V primeru klasifikacije uporabljamo drugačno funkcijo izgube, prečno entropijo
\[
  H = - \sum_i P_i \log_2 Q_i,
\]
kjer je $P$ prava (diskretna) porazdelitev cilje porazdelitve $Y$, torej $P_k =
1$ natanko tedaj, ko je trenutno opazovana vrednost enaka $k$, $Q$ pa je
porazdelitev, ki jo napove model, $Q_i = y_i^{L+1}$.
Potem je $E = - \log_2 y_k^{L+1}$, odvod pa
\[
  \frac{\partial E}{\partial \hat{y}_i}
  =
  \begin{cases}
	0 & i \ne k \\
	\inv{\hat{y}_k \ln 2} & i = k
  \end{cases}
\]
z enakim nadaljnjim delom.

\vprasanje{Kako izračunaš odvode za primer klasifikacijske nevronske mreže?}

Poznamo tudi konvolucijske nevronske mreže, kjer so nevroni v posamičnih plasteh
organizirani v matrike.
Namesto linearne obtežene vsote po vseh nevronih uporabimo konvolucijski filter,
tj.~matriko $c \times c$ uteži $w_{ij}$, s katero kombiniramo nevrone.
Če je plast $l-1$ dimenzije $x \times y$, je plast $l$ potem dimenzije $(x-c+1)
\times (y - c+1)$.
Pogosto želimo tudi zmanjšati dimenzije teh matrik, kar običajno naredimo z
max-akumulacijo, ki vzame največjo vrednost v neki izbrani podmatriki; ponavadi
velikosti $2 \times 2$.

\vprasanje{Opiši delovanje konvolucijskih nevronskih mrež.}

Poznamo tudi rekurenčne nevronske mreže, kjer dovolimo pojavitev zanke.
Za možnost računanja gradientov diskretiziramo čas in pri zanki za vhod
uporabimo izhod prejšnje iteracije; tako si lahko zapomnimo prejšnje podatke.
S tem smo pravzaprav dobili zaporedje enakih nevronskih mrež, ki ustrezajo
časovnim točkam.
Problem se pojavi pri učenju; napaka je sedaj funkcija časa
\[
  E = \inv{T} \sum_{t=1}^T L(y_t, \hat{y}_t),
\]
kjer je $L$ funkcija izgube.
Na podoben način kot prej lahko izračunamo odvod te napake, pri čemer dobimo
rekurzivno zvezo v $t$.

\vprasanje{Opiši rekurenčne nevronske mreže.}

% LocalWords:  perceptron softmax konvolucijske konvolucijski max-akumulacijo
% LocalWords:  konvolucijskih mrež rekurenčne diskretiziramo
